pdfoutput=1
% Options for packages loaded elsewhere
\PassOptionsToPackage{unicode}{hyperref}
\PassOptionsToPackage{hyphens}{url}
\pdfoutput=1
\documentclass[
  12pt,
]{article}
\usepackage{xcolor}
\usepackage{amsmath,amssymb}
\setcounter{secnumdepth}{-\maxdimen} % remove section numbering
\usepackage{iftex}
\ifPDFTeX
  \usepackage[T1]{fontenc}
  \usepackage{textcomp} % provide euro and other symbols
\else % if luatex or xetex
  \usepackage{unicode-math} % this also loads fontspec
  \defaultfontfeatures{Scale=MatchLowercase}
  \defaultfontfeatures[\rmfamily]{Ligatures=TeX,Scale=1}
\fi
\usepackage{lmodern}
\ifPDFTeX\else
  % xetex/luatex font selection
\fi
% Use upquote if available, for straight quotes in verbatim environments
\IfFileExists{upquote.sty}{\usepackage{upquote}}{}
\IfFileExists{microtype.sty}{% use microtype if available
  \usepackage[]{microtype}
  \UseMicrotypeSet[protrusion]{basicmath} % disable protrusion for tt fonts
}{}
\makeatletter
\@ifundefined{KOMAClassName}{% if non-KOMA class
  \IfFileExists{parskip.sty}{%
    \usepackage{parskip}
  }{% else
    \setlength{\parindent}{0pt}
    \setlength{\parskip}{6pt plus 2pt minus 1pt}}
}{% if KOMA class
  \KOMAoptions{parskip=half}}
\makeatother
\usepackage{longtable,booktabs,array}
\usepackage{calc} % for calculating minipage widths
% Correct order of tables after \paragraph or \subparagraph
\usepackage{etoolbox}
\makeatletter
\patchcmd\longtable{\par}{\if@noskipsec\mbox{}\fi\par}{}{}
\makeatother
% Allow footnotes in longtable head/foot
\IfFileExists{footnotehyper.sty}{\usepackage{footnotehyper}}{\usepackage{footnote}}
\makesavenoteenv{longtable}
\setlength{\emergencystretch}{3em} % prevent overfull lines
\providecommand{\tightlist}{%
  \setlength{\itemsep}{0pt}\setlength{\parskip}{0pt}}
\usepackage{bookmark}
\IfFileExists{xurl.sty}{\usepackage{xurl}}{} % add URL line breaks if available
\urlstyle{same}
\hypersetup{
  hidelinks,
  pdfcreator={LaTeX via pandoc}}

\author{SCX}
\date{2026}

\begin{document}

\section{Connections Between SCX Self-Evolution and Known
Theories}\label{connections-between-scx-self-evolution-and-known-theories}

\begin{quote}
\textbf{Part of the SCX Self-Evolution Theory Series} \textbf{Status}:
Formal comparison \textbar{} \textbf{Audit}: Pre-review
\textbf{Prerequisites}: THEOREMS\_UNIFIED.md, self-evolution definitions
(Files 1-7)
\end{quote}

\begin{center}\rule{0.5\linewidth}{0.5pt}\end{center}

\subsection{Table of Contents}\label{table-of-contents}

\begin{enumerate}
\def\labelenumi{\arabic{enumi}.}
\tightlist
\item
  \hyperref[1-alphazero-self-play]{AlphaZero Self-Play}
\item
  \hyperref[2-bayesian-optimization]{Bayesian Optimization}
\item
  \hyperref[3-active-learning]{Active Learning}
\item
  \hyperref[4-solomonoff-induction]{Solomonoff Induction}
\item
  \hyperref[5-comprehensive-comparison-table]{Comprehensive Comparison
  Table}
\item
  \hyperref[6-synthesis-what-is-and-is-not-new]{Synthesis: What Is and
  Is Not New}
\end{enumerate}

\begin{center}\rule{0.5\linewidth}{0.5pt}\end{center}

\subsection{1. AlphaZero Self-Play}\label{alphazero-self-play}

\subsubsection{1.1 Overview}\label{overview}

AlphaZero (Silver et al., 2017, 2018) learns to play board games through
self-play reinforcement learning. The agent plays games against itself,
using Monte Carlo Tree Search (MCTS) to generate training data, which is
then used to update a deep neural network. The network has two heads: a
policy head \(p(s)\) predicting move probabilities and a value head
\(V(s)\) predicting the expected outcome from state \(s\).

The self-play loop is:
\[\pi_{\theta_{\text{old}}} \xrightarrow{\text{MCTS}} \{\text{self-play games}\} \xrightarrow{\text{training}} \pi_\theta \xrightarrow{\text{evaluation}} \pi_{\theta_{\text{new}}}\]

\subsubsection{1.2 Formal Mapping}\label{formal-mapping}

The SCX self-evolution loop maps to AlphaZero as follows:

\begin{longtable}[]{@{}
  >{\raggedright\arraybackslash}p{(\linewidth - 4\tabcolsep) * \real{0.3279}}
  >{\raggedright\arraybackslash}p{(\linewidth - 4\tabcolsep) * \real{0.4918}}
  >{\raggedright\arraybackslash}p{(\linewidth - 4\tabcolsep) * \real{0.1803}}@{}}
\toprule\noalign{}
\begin{minipage}[b]{\linewidth}\raggedright
AlphaZero Component
\end{minipage} & \begin{minipage}[b]{\linewidth}\raggedright
SCX Self-Evolution Component
\end{minipage} & \begin{minipage}[b]{\linewidth}\raggedright
Rationale
\end{minipage} \\
\midrule\noalign{}
\endhead
\bottomrule\noalign{}
\endlastfoot
Policy network \(\pi_\theta\) & Gatekeeper \(S_t\) & Both select
actions: \(\pi_\theta\) selects moves, \(S_t\) selects data for
validation \\
Self-play games & NEP student \(f_{\theta_t}\) predictions & Both
generate training targets through interaction with a fixed
environment \\
MCTS search & SCX state-conditioned aggregation
\(\hat{R}_m(s), \text{SCX}_m(s)\) & Both refine raw evaluations through
structured computation \\
Value network \(V(s)\) & SCX state value function
\(V(s) = \mathbb{E}[\text{improvement} \mid s]\) & Both estimate the
long-term value of states/decisions \\
Replay buffer & Memory bank \(M_t\) & Both store experience for future
training \\
Evaluation against previous version & Lyapunov descent check
\(\Phi(S_t) < \Phi(S_{t-1})\) & Both compare new policy against old to
ensure monotonic improvement \\
Opponent (fixed rules) & True physical law \(f^*\) + current data
\(\mathcal{D}_t\) & Both provide the fixed, non-learnable environment \\
Training loss
\(L(\theta) = L_{\text{policy}} + L_{\text{value}} + L_{\text{reg}}\) &
Lyapunov function \(\Phi(S_t, M_t, f_{\theta_t})\) & Both provide the
scalar signal for optimization \\
\end{longtable}

\subsubsection{1.3 Key Similarities}\label{key-similarities}

\begin{enumerate}
\def\labelenumi{\arabic{enumi}.}
\item
  \textbf{Bootstrapping}: Both systems generate their own training data.
  AlphaZero's training data comes from MCTS self-play; SCX's training
  data comes from the NEP student's self-generated configurations with
  gatekeeper-selected validation.
\item
  \textbf{Monotonic improvement}: AlphaZero uses policy iteration, which
  guarantees non-decreasing performance against the previous version.
  SCX uses Lyapunov descent, which guarantees
  \(\Phi(S_t) \leq \Phi(S_{t-1})\).
\item
  \textbf{Memory buffer}: Both store past experience. AlphaZero's replay
  buffer prevents catastrophic forgetting; SCX's memory bank \(M_t\)
  preserves previously validated data.
\item
  \textbf{Exploration-exploitation}: AlphaZero's MCTS balances
  exploration (via Dirichlet noise in prior probabilities) with
  exploitation (via value-guided search). SCX's gatekeeper balances
  exploration (validating uncertain states) with exploitation (trusting
  states with high SCX reliability scores).
\end{enumerate}

\subsubsection{1.4 Key Differences}\label{key-differences}

\begin{enumerate}
\def\labelenumi{\arabic{enumi}.}
\item
  \textbf{Data generation mechanism}: AlphaZero generates
  \textbf{synthetic} games through self-play. The environment (game
  rules) is completely known and simulable. SCX generates data through
  NEP \textbf{simulations} of physical systems. The environment
  (physical law \(f^*\)) is \textbf{unknown} and only partially
  observable through noisy labels.
\item
  \textbf{Feedback delay}: AlphaZero receives immediate game outcomes
  (win/loss/draw) at the end of each game. The value target is
  unambiguous. SCX's feedback from NEP validation is \textbf{delayed}
  and \textbf{partial}: only a subset of samples are validated, and the
  validation itself is noisy (DFT approximations, experimental error).
\item
  \textbf{Convergence mechanism}: AlphaZero uses \textbf{policy
  iteration} --- solving a series of increasingly accurate supervised
  learning problems against a fixed opponent. SCX uses
  \textbf{memory-augmented Bayesian update} --- the gatekeeper is
  updated based on accumulated evidence in \(M_t\), and the NEP student
  is updated based on the full memory bank.
\item
  \textbf{Distribution shift}: AlphaZero's data distribution shifts as
  the policy improves, but the game dynamics remain fixed. SCX's
  distribution shifts \textbf{doubly}: the gatekeeper changes which data
  are selected for validation, and the NEP student changes the space of
  accessible configurations.
\item
  \textbf{Final state}: AlphaZero converges to a superhuman policy (in
  principle, the Nash equilibrium of the game if the network has
  sufficient capacity). SCX converges to a self-consistent fixed point
  (Theorem SE-2) that may be \textbf{suboptimal} due to the finite
  configuration space and the fundamental unidentifiability of noise
  (Theorem 3).
\end{enumerate}

\subsubsection{1.5 Convergence Comparison}\label{convergence-comparison}

\textbf{AlphaZero (policy iteration)}:
\[\text{Performance}(\pi_{t+1}) \geq \text{Performance}(\pi_t)\]
Convergence to optimal policy \(\pi^*\) under sufficient capacity and
exploration (proven for tabular case; empirically observed for neural
networks).

\textbf{SCX (Lyapunov descent)}:
\[\Phi(S_{t+1}, M_{t+1}, f_{\theta_{t+1}}) \leq \Phi(S_t, M_t, f_{\theta_t})\]
Convergence to fixed point \(q_{T^*}\) in finite time (Theorem SE-2).
The fixed point is self-consistent but not necessarily optimal in a
global sense.

\textbf{Comparison}: Both guarantee monotonic improvement. AlphaZero's
guarantee is stronger (convergence to optimal play under certain
conditions) but relies on a fully known environment. SCX's guarantee is
weaker (convergence to a self-consistent fixed point) but operates in an
environment with unknown ground truth.

\begin{center}\rule{0.5\linewidth}{0.5pt}\end{center}

\subsection{2. Bayesian Optimization}\label{bayesian-optimization}

\subsubsection{2.1 Overview}\label{overview-1}

Bayesian optimization (BO) (Mockus, 1975; Jones et al., 1998) optimizes
a black-box function \(f: \mathcal{X} \to \mathbb{R}\) by building a
probabilistic surrogate model \(g\) (typically a Gaussian process) and
using an acquisition function \(\alpha\) to select the next evaluation
point:

\[x_{t+1} = \arg\max_{x \in \mathcal{X}} \alpha(x \mid g_{1:t})\]

where \(\alpha\) balances exploration (points with high uncertainty) and
exploitation (points with high predicted value).

\subsubsection{2.2 Formal Mapping}\label{formal-mapping-1}

\begin{longtable}[]{@{}
  >{\raggedright\arraybackslash}p{(\linewidth - 4\tabcolsep) * \real{0.2407}}
  >{\raggedright\arraybackslash}p{(\linewidth - 4\tabcolsep) * \real{0.5556}}
  >{\raggedright\arraybackslash}p{(\linewidth - 4\tabcolsep) * \real{0.2037}}@{}}
\toprule\noalign{}
\begin{minipage}[b]{\linewidth}\raggedright
BO Component
\end{minipage} & \begin{minipage}[b]{\linewidth}\raggedright
SCX Self-Evolution Component
\end{minipage} & \begin{minipage}[b]{\linewidth}\raggedright
Rationale
\end{minipage} \\
\midrule\noalign{}
\endhead
\bottomrule\noalign{}
\endlastfoot
Black-box function \(f\) & True physical law \(f^*\) & Both are unknown,
expensive to evaluate \\
Surrogate model \(g\) & NEP student \(f_\theta\) & Both approximate the
unknown ground truth \\
Acquisition function \(\alpha\) & Gatekeeper \(S_t\) & Both guide the
selection of the next evaluation point \\
Evaluation budget & NEP validation budget & Both have limited
evaluations (BO: function evaluations; SCX: DFT/experiment budget) \\
Posterior \(p(f \mid \mathcal{D}_{1:t})\) & Memory bank \(M_t\) +
student \(f_{\theta_t}\) & Both represent accumulated knowledge about
the unknown function \\
Regret \(r_t = f(x^*) - f(x_t)\) & Lyapunov gap
\(\Phi(q_t) - \Phi_{\text{opt}}\) & Both measure suboptimality of the
current state \\
\end{longtable}

\subsubsection{2.3 Formal Comparison: Acquisition
Functions}\label{formal-comparison-acquisition-functions}

\textbf{BO's Expected Improvement (EI)}:
\[\alpha_{\text{EI}}(x) = \mathbb{E}\left[\max(0, f(x) - f_{\text{best}}) \mid \mathcal{D}_{1:t}\right]\]

\textbf{SCX's State Data Value}:
\[V(s) = \underbrace{\hat{R}(s)}_{\text{expert disagreement}} \cdot \underbrace{\rho(s)}_{\text{state proportion}} \cdot \underbrace{(1 - C(s))}_{\text{uncertainty}}\]

where \(C(s) = \frac{1}{M}\sum_m \mathbf{1}\{\ell(f_m(x), y) > \tau\}\)
is the consensus score (Theorem 1).

Both EI and SCX's \(V(s)\) balance: - \textbf{Exploitation} (high
predicted improvement / high expert disagreement \(\hat{R}(s)\)) -
\textbf{Exploration} (high posterior uncertainty / low consensus
\(C(s)\))

\textbf{Key difference}: EI operates on a continuous input space and is
formalized as an expectation under a Gaussian process posterior. SCX's
\(V(s)\) operates on a discrete state space (partition of
\(\mathcal{X}\) into \(K_S\) states) and uses frequentist estimates of
within-state disagreement.

\subsubsection{2.4 Convergence Rates}\label{convergence-rates}

\textbf{BO convergence} (Srinivas et al., 2010): For a GP with kernel
\(k\), the cumulative regret is bounded by:

\[R_T = \sum_{t=1}^T r_t \leq \sqrt{T \cdot \gamma_T \cdot C}\]

where \(\gamma_T\) is the maximum information gain after \(T\)
evaluations. For the squared exponential kernel,
\(\gamma_T = O((\log T)^{d+1})\), giving
\(R_T = O(\sqrt{T (\log T)^{d+1}})\).

\textbf{SCX convergence} (Theorem SE-1, SE-2): For the Lyapunov function
\(\Phi\), the gap decays as:

\[\Phi(q_t) - \Phi_{\text{opt}} \leq \Phi(q_0) \cdot \exp(-\lambda t)\]

under the contraction condition, or in worst case:

\[T^* \leq \frac{\Phi_0}{\varepsilon_{\text{mach}}}\]

\textbf{Comparison}: BO provides cumulative regret bounds that depend on
the information-theoretic complexity of the function class. SCX provides
a finite-time termination bound that depends on the Lyapunov descent
rate and machine precision. BO's bounds are more informative for
sub-exponential convergence; SCX's guarantee is stronger (exact
termination).

\subsubsection{2.5 Key Similarities}\label{key-similarities-1}

\begin{enumerate}
\def\labelenumi{\arabic{enumi}.}
\tightlist
\item
  \textbf{Uncertainty-guided exploration}: Both systems use uncertainty
  estimates to allocate evaluation resources where they are most needed.
\item
  \textbf{Surrogate-based decision making}: Both maintain an internal
  model (GP / NEP student) that guides data acquisition.
\item
  \textbf{Balanced exploration-exploitation}: Both use explicit
  mechanisms to balance exploring unknown regions against exploiting
  known good regions.
\item
  \textbf{Sequential design}: Both are iterative, adding one (or a
  batch) of evaluation points per iteration.
\end{enumerate}

\subsubsection{2.6 Key Differences}\label{key-differences-1}

\begin{enumerate}
\def\labelenumi{\arabic{enumi}.}
\item
  \textbf{Target distribution}: BO optimizes a \textbf{static} black-box
  function. The objective \(f\) is fixed across iterations. SCX faces an
  \textbf{evolving} target distribution: as the NEP student improves,
  the set of ``interesting'' configurations changes.
\item
  \textbf{Batch selection}: SCX typically validates an entire batch of
  samples (all samples in a selected state) before updating the
  gatekeeper. BO typically evaluates one point at a time (though batch
  BO variants exist).
\item
  \textbf{State structure}: SCX leverages the state partition \(S\) to
  aggregate information across similar inputs. BO treats each input
  point independently (though GP kernels induce similarity structure).
\item
  \textbf{Multiple objectives}: SCX's gatekeeper manages multiple
  criteria (noise detection, expert disagreement, data diversity). BO
  typically optimizes a single scalar objective (though multi-objective
  BO variants exist).
\item
  \textbf{Cost structure}: BO typically assumes homogeneous evaluation
  costs. SCX's validation cost varies across states (DFT calculations
  for some configurations are more expensive than others).
\end{enumerate}

\begin{center}\rule{0.5\linewidth}{0.5pt}\end{center}

\subsection{3. Active Learning}\label{active-learning}

\subsubsection{3.1 Overview}\label{overview-2}

Active learning (AL) (Settles, 2009) aims to reduce labeling effort by
selectively querying labels for the most informative unlabeled examples.
Given a labeled pool \(\mathcal{L}\), an unlabeled pool \(\mathcal{U}\),
and a model \(h\), the query strategy \(Q(x \mid h)\) selects the next
sample(s) to label.

\subsubsection{3.2 Formal Mapping}\label{formal-mapping-2}

\begin{longtable}[]{@{}
  >{\raggedright\arraybackslash}p{(\linewidth - 4\tabcolsep) * \real{0.2407}}
  >{\raggedright\arraybackslash}p{(\linewidth - 4\tabcolsep) * \real{0.5556}}
  >{\raggedright\arraybackslash}p{(\linewidth - 4\tabcolsep) * \real{0.2037}}@{}}
\toprule\noalign{}
\begin{minipage}[b]{\linewidth}\raggedright
AL Component
\end{minipage} & \begin{minipage}[b]{\linewidth}\raggedright
SCX Self-Evolution Component
\end{minipage} & \begin{minipage}[b]{\linewidth}\raggedright
Rationale
\end{minipage} \\
\midrule\noalign{}
\endhead
\bottomrule\noalign{}
\endlastfoot
Model \(h\) & NEP student \(f_\theta\) & Both are the predictive model
being improved \\
Oracle \(O\) (labeler) & NEP simulator + DFT validation & Both provide
ground-truth labels \\
Query strategy \(Q(x \mid h)\) & Gatekeeper \(S_t\) + state-value
scoring & Both select which samples to label/validate \\
Labeled pool \(\mathcal{L}\) & Memory bank \(M_t\) & Both store
labeled/validated data \\
Unlabeled pool \(\mathcal{U}\) & Candidate pool \(\mathcal{C}_t\) & Both
contain unlabeled samples \\
Labeling budget & NEP validation budget & Both have finite labeling
capacity \\
\end{longtable}

\subsubsection{3.3 Comparison of Query
Strategies}\label{comparison-of-query-strategies}

\textbf{Standard AL strategies}:

\begin{enumerate}
\def\labelenumi{\arabic{enumi}.}
\item
  \textbf{Uncertainty sampling}: \(Q(x) = 1 - \max_y P(y \mid x)\),
  selects samples where the model is most uncertain.
\item
  \textbf{Query-by-committee (QBC)}:
  \(Q(x) = \frac{1}{M}\sum_{m=1}^M \mathbf{1}\{h_m(x) \neq \bar{h}(x)\}\),
  selects samples with maximum disagreement among an ensemble of models.
\item
  \textbf{Expected model change}:
  \(Q(x) = \mathbb{E}_{y \sim P(y|x)}[\|\nabla L(\theta \cup (x,y)) - \nabla L(\theta)\|]\),
  selects samples expected to maximally change the model.
\end{enumerate}

\textbf{SCX's approach (state-certified active learning)}:

\[a^*(s) = \arg\max_{a \in \mathcal{A}} U(a, s)\]

where
\(\mathcal{A} = \{\text{validate}, \text{discard}, \text{defer}, \ldots\}\)
and \(U(a, s)\) is the expected utility of action \(a\) for state \(s\).
The utility is computed from state-level statistics:

\[U(\text{validate}, s) = \underbrace{\hat{R}(s)}_{\text{expert disagreement}} \cdot \underbrace{(1 - \text{SCX}(s))}_{\text{reliability gap}} \cdot \underbrace{\rho(s)}_{\text{state size}}\]

\subsubsection{3.4 Comparison of Query
Strategies}\label{comparison-of-query-strategies-1}

\begin{longtable}[]{@{}
  >{\raggedright\arraybackslash}p{(\linewidth - 6\tabcolsep) * \real{0.1852}}
  >{\raggedright\arraybackslash}p{(\linewidth - 6\tabcolsep) * \real{0.2778}}
  >{\raggedright\arraybackslash}p{(\linewidth - 6\tabcolsep) * \real{0.2407}}
  >{\raggedright\arraybackslash}p{(\linewidth - 6\tabcolsep) * \real{0.2963}}@{}}
\toprule\noalign{}
\begin{minipage}[b]{\linewidth}\raggedright
Strategy
\end{minipage} & \begin{minipage}[b]{\linewidth}\raggedright
AL Counterpart
\end{minipage} & \begin{minipage}[b]{\linewidth}\raggedright
SCX Version
\end{minipage} & \begin{minipage}[b]{\linewidth}\raggedright
Key Difference
\end{minipage} \\
\midrule\noalign{}
\endhead
\bottomrule\noalign{}
\endlastfoot
Uncertainty sampling & \(1 - \max_y P(y \mid x)\) & \(1 - C(s)\) where
\(C(s)\) is state consensus & SCX uses multi-expert consensus instead of
single-model uncertainty \\
Query-by-committee & Vote entropy among ensemble & SCX reliability gap
\(\text{SCX}(s)\) & SCX uses state-conditioned expert reliability, not
raw vote entropy \\
Expected model change &
\(\|\nabla L(\theta \cup (x,y)) - \nabla L(\theta)\|\) &
\(V(s) = \hat{R}(s) \cdot \rho(s) \cdot (1 - C(s))\) & SCX's value
function operates at state level, not point level \\
Density-weighted & \(\frac{1}{n}\sum_{i=1}^n \text{sim}(x, x_i)\) &
\(\rho(s)\) (state proportion) & Both weight by representativeness \\
Diversity sampling & Coreset selection & Discard redundant states via
SCX-Compress & Both aim for coverage \\
\end{longtable}

\subsubsection{3.5 SCX's State-Certified Active
Learning}\label{scxs-state-certified-active-learning}

The key innovation of SCX's approach is the \textbf{state-level
certification}. Unlike point-wise active learning, which treats each
sample independently:

\begin{enumerate}
\def\labelenumi{\arabic{enumi}.}
\item
  \textbf{State aggregation}: SCX estimates statistics at the state
  level, pooling information across all samples in state \(s\). This
  provides more robust estimates when individual labels are noisy.
\item
  \textbf{Certified actions}: Each action \(a \in \mathcal{A}\) is
  certified by state-level guarantees. For example, ``validate state
  \(s\)'' is justified when \(\text{SCX}(s)\) is below threshold,
  indicating that the multi-expert consensus is unreliable for this
  state.
\item
  \textbf{Cascading decisions}: SCX makes decisions hierarchically:
  first at the state level (which state to focus on), then at the sample
  level (which samples within the state to validate).
\end{enumerate}

\subsubsection{3.6 Sample Complexity
Comparison}\label{sample-complexity-comparison}

\textbf{Standard AL sample complexity} (Balcan et al., 2009): Under the
disagreement coefficient \(\theta\), active learning achieves label
complexity:

\[N_{\text{AL}}(\varepsilon) = O\left(\theta \cdot d \cdot \log\frac{1}{\varepsilon}\right)\]

where \(d\) is the VC dimension of the hypothesis class.

\textbf{SCX sample complexity} (derived from Theorem 1 and Proposition
2): The number of validation labels needed per state \(s\) to achieve
reliability \(\text{SCX}_m(s) > 1 - \delta\) is:

\[n_{\text{SCX}}(s, \delta) = O\left(\frac{1}{\Delta_s^2} \log\frac{1}{\delta}\right)\]

where \(\Delta_s\) is the state-level separation gap (Theorem 1).
Aggregating over \(K_S\) states:

\[N_{\text{SCX}}(\varepsilon) = O\left(K_S \cdot \frac{1}{\Delta_{\min}^2} \log\frac{1}{\varepsilon}\right)\]

where \(\Delta_{\min} = \min_s \Delta_s\).

\textbf{Comparison}: Standard AL avoids the \(K_S\) multiplicative
factor but requires the disagreement coefficient \(\theta\) to be small.
SCX's factor of \(K_S\) is the cost of operating without a single
discriminative model --- SCX uses multiple experts and must characterize
each state independently.

\begin{center}\rule{0.5\linewidth}{0.5pt}\end{center}

\subsection{4. Solomonoff Induction}\label{solomonoff-induction}

\subsubsection{4.1 Overview}\label{overview-3}

Solomonoff induction (Solomonoff, 1964) is a theoretical framework for
inductive inference. Given a sequence of observations, the posterior
probability of a continuation is:

\[P(x_{n+1} \mid x_1, \ldots, x_n) = \frac{\sum_{p: p(x_1,\ldots,x_n) = x_1,\ldots,x_n} 2^{-|p|} \cdot P(x_{n+1} \mid p)}{\sum_{p: p(x_1,\ldots,x_n) = x_1,\ldots,x_n} 2^{-|p|}}\]

where \(p\) ranges over all programs for a universal Turing machine,
\(|p|\) is program length, and \(P(x_{n+1} \mid p)\) is the continuation
probability under program \(p\). The universal prior
\(M(x) = \sum_p 2^{-|p|} \cdot \mathbf{1}\{p(\varepsilon) = x\}\)
assigns higher probability to strings generated by shorter programs
(Occam's razor).

\subsubsection{4.2 Formal Analogy}\label{formal-analogy}

\textbf{Claim (SE-A1)}: The limiting SCX gatekeeper
\(S_\infty(x) = \lim_{t\to\infty} S_t(x)\) approximates a Solomonoff
predictor for the question ``Is this label correct?''

The analogy proceeds as follows:

\begin{longtable}[]{@{}
  >{\raggedright\arraybackslash}p{(\linewidth - 4\tabcolsep) * \real{0.2683}}
  >{\raggedright\arraybackslash}p{(\linewidth - 4\tabcolsep) * \real{0.4634}}
  >{\raggedright\arraybackslash}p{(\linewidth - 4\tabcolsep) * \real{0.2683}}@{}}
\toprule\noalign{}
\begin{minipage}[b]{\linewidth}\raggedright
Solomonoff
\end{minipage} & \begin{minipage}[b]{\linewidth}\raggedright
SCX Self-Evolution
\end{minipage} & \begin{minipage}[b]{\linewidth}\raggedright
Rationale
\end{minipage} \\
\midrule\noalign{}
\endhead
\bottomrule\noalign{}
\endlastfoot
Universal prior
\(M(x) = \sum_p 2^{-|p|}\mathbf{1}\{p(\varepsilon) = x\}\) & State
partition \(\mathcal{S}\) with prior \(\rho_s = \mathbb{P}(X \in s)\) &
Both define prior beliefs over hypotheses \\
Programs \(p\) & Expert models \(\{f_m\}\) & Both are candidate
explanations of observed data \\
Program length \(|p|\) (Occam penalty) & Expert complexity (implicit in
training data requirement) & Both penalize complexity \\
Likelihood \(P(\text{data} \mid p)\) & Consensus score
\(C(x) = \frac{1}{M}\sum_m \mathbf{1}\{\ell(f_m(x), y) > \tau\}\) & Both
measure how well the hypothesis explains the observation \\
Posterior \(P(x_{n+1} \mid \text{data})\) & Gatekeeper score \(S_t(x)\)
& Both represent the current belief about an unseen property \\
Accumulated evidence & Memory bank \(M_t\) & Both store the history of
observations \\
Convergence as \(n \to \infty\) & Convergence as \(t \to \infty\)
(Theorem SE-2) & Both converge to fixed points given sufficient
evidence \\
\end{longtable}

\subsubsection{4.3 Formal Correspondence}\label{formal-correspondence}

The Solomonoff posterior for the proposition ``\(x\) is clean'' (i.e.,
\(y = f^*(x)\)) is:

\[P(\text{clean} \mid \mathcal{D}) = \frac{\sum_p 2^{-|p|} \cdot P(\mathcal{D} \mid p) \cdot \mathbf{1}\{p \text{ predicts clean at } x\}}{\sum_p 2^{-|p|} \cdot P(\mathcal{D} \mid p)}\]

The SCX gatekeeper's score for \(x\) at state \(s(x)\) is:

\[S_t(x) = \frac{1}{M}\sum_{m=1}^M \underbrace{\text{SCX}_m(s(x))}_{\text{expert reliability in state } s} \cdot \underbrace{\mathbf{1}\{\ell(f_m(x), y) < \tau\}}_{\text{expert } m \text{ agrees with label}}\]

The correspondence is: - \textbf{Hypothesis class}: Solomonoff sums over
all computable functions; SCX sums over the finite set of \(M\) experts.
- \textbf{Prior}: Solomonoff uses the universal prior \(2^{-|p|}\); SCX
uses state proportions \(\rho_s\) (experts are implicitly weighted by
their state-conditioned reliability). - \textbf{Convergence}: Solomonoff
converges to the truth with probability 1 for any computable environment
(Hutter, 2005). SCX converges to a fixed point (Theorem SE-2) but may
not converge to the truth (Theorem 3).

\subsubsection{4.4 Key Limitations of the
Analogy}\label{key-limitations-of-the-analogy}

\begin{enumerate}
\def\labelenumi{\arabic{enumi}.}
\item
  \textbf{Hypothesis class restriction}: Solomonoff induction sums over
  \textbf{all} computable functions, making it a universal learner
  (incomputable but theoretically optimal). SCX sums over only \(M\)
  experts, which is a \textbf{severely restricted} hypothesis class. The
  SCX gatekeeper cannot represent a belief that none of the \(M\)
  experts captures.

  \[S_\infty(x) \approx \text{Solomonoff predictor on a restricted hypothesis class}\]
\item
  \textbf{Computability}: Solomonoff induction is incomputable (the sum
  over all programs requires solving the halting problem). SCX's
  gatekeeper update is efficiently computable (polynomial in \(M\),
  \(K_S\), and \(N\)).
\item
  \textbf{No formal Occam prior}: SCX does not explicitly implement
  Occam's razor. Simpler state structures are preferred
  \textbf{implicitly} (k-means with few clusters is more stable;
  Proposition 6 favors reproducible partitions), but there is no formal
  simplicity penalty.
\item
  \textbf{No complete theory of convergence to truth}: Theorem 3 shows
  that SCX cannot always distinguish noise from difficulty, even in the
  infinite-data limit. Solomonoff induction does not have this
  limitation (any computable pattern is eventually learned).
\end{enumerate}

\subsubsection{4.5 Occam's Razor in SCX}\label{occams-razor-in-scx}

While SCX does not explicitly implement Occam's razor, simpler state
structures are implicitly preferred through:

\begin{itemize}
\item
  \textbf{k-means stability} (Proposition 6): Partitions with few,
  well-separated states have higher stability scores \(S(\Phi, K_S)\).
  The stability diagnostic implicitly favors simpler state structures.
\item
  \textbf{Cluster count selection}: Using the stability score to select
  \(K_S\) (choosing the smallest \(K_S\) with acceptable stability)
  implements a form of model selection that penalizes complexity.
\item
  \textbf{Bootstrap consistency}: The bootstrap diagnostic (Proposition
  6) detects when additional states are spurious (low ARI across
  resamples), discouraging overfitting in state discovery.
\end{itemize}

This implicit Occam bias is weaker than Solomonoff's explicit prior
\(2^{-|p|}\) but serves a similar function: it prevents the gatekeeper
from overfitting to noise in the expert consensus signals.

\begin{center}\rule{0.5\linewidth}{0.5pt}\end{center}

\subsection{5. Comprehensive Comparison
Table}\label{comprehensive-comparison-table}

\subsubsection{5.1 Unified Comparison}\label{unified-comparison}

\begin{longtable}[]{@{}
  >{\raggedright\arraybackslash}p{(\linewidth - 10\tabcolsep) * \real{0.1100}}
  >{\raggedright\arraybackslash}p{(\linewidth - 10\tabcolsep) * \real{0.1100}}
  >{\raggedright\arraybackslash}p{(\linewidth - 10\tabcolsep) * \real{0.2100}}
  >{\raggedright\arraybackslash}p{(\linewidth - 10\tabcolsep) * \real{0.1700}}
  >{\raggedright\arraybackslash}p{(\linewidth - 10\tabcolsep) * \real{0.2100}}
  >{\raggedright\arraybackslash}p{(\linewidth - 10\tabcolsep) * \real{0.1900}}@{}}
\toprule\noalign{}
\begin{minipage}[b]{\linewidth}\raggedright
Dimension
\end{minipage} & \begin{minipage}[b]{\linewidth}\raggedright
AlphaZero
\end{minipage} & \begin{minipage}[b]{\linewidth}\raggedright
Bayesian Optimization
\end{minipage} & \begin{minipage}[b]{\linewidth}\raggedright
Active Learning
\end{minipage} & \begin{minipage}[b]{\linewidth}\raggedright
Solomonoff Induction
\end{minipage} & \begin{minipage}[b]{\linewidth}\raggedright
SCX Self-Evolution
\end{minipage} \\
\midrule\noalign{}
\endhead
\bottomrule\noalign{}
\endlastfoot
\textbf{Core Mechanism} & Policy iteration + MCTS + neural network &
Surrogate model + acquisition function & Query strategy + oracle
labeling & Universal prior + Bayesian update & Lyapunov descent + memory
bank + gatekeeper \\
\textbf{Exploration Strategy} & Dirichlet noise + MCTS visit counts &
Posterior variance (GP) & Uncertainty/disagreement/diversity sampling &
Universal prior over all programs & State value function \(V(s)\) +
consensus uncertainty \\
\textbf{Convergence Guarantee} & Monotonic policy improvement; optimal
for tabular case & \(\sqrt{T \cdot \gamma_T}\) cumulative regret; no
exact termination & \(\tilde{O}(\theta \cdot d)\) label complexity;
VC-based & Convergence to truth with prob 1 for computable environments
& Finite-time fixed point (Theorem SE-2); Lyapunov descent (Theorem
SE-1) \\
\textbf{Regret Bound} & \(V^* - V^{\pi_t} \leq \text{unknown}\)
(empirically fast) & \(R_T = O(\sqrt{T \cdot \gamma_T})\) (GP) & Label
complexity \(O(\theta \cdot d \cdot \log(1/\varepsilon))\) & Optimal but
incomputable &
\(\Phi(q_t) - \Phi_{\text{opt}} \leq \Phi(q_0) e^{-\lambda t}\) (under
contraction) \\
\textbf{Computational Cost} & \(O(N_{\text{sim}} \cdot d_{\text{net}})\)
per iteration & \(O(T \cdot (n^3 + dn))\) for GP & \(O(n \cdot d)\) per
query & Incomputable & \(O(M \cdot K_S \cdot N + K_S \cdot d_\phi)\) per
iteration \\
\textbf{Key Difference from SCX} & Synthetic data; known environment;
stronger convergence & Static black-box; point-wise; single objective &
Point-wise queries; single model; no state aggregation & Universal class
but incomputable; complete convergence & Real data; unknown \(f^*\);
state-level; multi-objective; finite compute \\
\textbf{Environment} & Fully known (game rules) & Partially known (GP
prior) & Partially known (labeled pool) & Unknown but computable &
Unknown physical law (uncomputable in general) \\
\end{longtable}

\subsubsection{5.2 Detailed Metric
Comparison}\label{detailed-metric-comparison}

\begin{longtable}[]{@{}
  >{\raggedright\arraybackslash}p{(\linewidth - 10\tabcolsep) * \real{0.2051}}
  >{\raggedright\arraybackslash}p{(\linewidth - 10\tabcolsep) * \real{0.2821}}
  >{\raggedright\arraybackslash}p{(\linewidth - 10\tabcolsep) * \real{0.1282}}
  >{\raggedright\arraybackslash}p{(\linewidth - 10\tabcolsep) * \real{0.1282}}
  >{\raggedright\arraybackslash}p{(\linewidth - 10\tabcolsep) * \real{0.1282}}
  >{\raggedright\arraybackslash}p{(\linewidth - 10\tabcolsep) * \real{0.1282}}@{}}
\toprule\noalign{}
\begin{minipage}[b]{\linewidth}\raggedright
Metric
\end{minipage} & \begin{minipage}[b]{\linewidth}\raggedright
AlphaZero
\end{minipage} & \begin{minipage}[b]{\linewidth}\raggedright
BO
\end{minipage} & \begin{minipage}[b]{\linewidth}\raggedright
AL
\end{minipage} & \begin{minipage}[b]{\linewidth}\raggedright
SI
\end{minipage} & \begin{minipage}[b]{\linewidth}\raggedright
SCX
\end{minipage} \\
\midrule\noalign{}
\endhead
\bottomrule\noalign{}
\endlastfoot
\textbf{Per-iteration data cost} & Free (synthetic) & \(O(1)\) function
evaluation & \(O(b)\) labels per batch & None & \(O(b)\) NEP
validations \\
\textbf{Scalability (\(X\) dimension)} & High (neural nets) & Low (GP:
\(O(n^3)\)) & Medium & None (incomputable) & Medium (k-means:
\(O(N d_\phi K_S)\)) \\
\textbf{Theoretical maturity} & Empirical (except tabular) & Strong (GP
regret) & Strong (disagreement coeff.) & Complete (incomputable) &
Partial (this work) \\
\textbf{Applicability to physical sciences} & Not designed (game domain)
& High (materials, chemistry) & Medium (active labeling) & None
(incomputable) & High (designed for NEP) \\
\textbf{Handles label noise} & No (synthetic data) & No (GP assumes
clean) & Partially (oracle assumed clean) & Yes (Bayesian) & Yes
(primary design goal) \\
\textbf{Uncertainty quantification} & Value head variance & GP posterior
& Model confidence & Full posterior & Multi-expert consensus \\
\textbf{Memory of past data} & Replay buffer (bounded) & All past
evaluations (\(O(n)\)) & Labeled pool & All past data & Memory bank
\(M_t\) \\
\textbf{Formal convergence proof} & Partial (tabular) & Yes (GP regret)
& Yes (label complexity) & Yes (incomputable) & Yes (SE-1, SE-2) \\
\end{longtable}

\begin{center}\rule{0.5\linewidth}{0.5pt}\end{center}

\subsection{6. Synthesis: What Is and Is Not
New}\label{synthesis-what-is-and-is-not-new}

\subsubsection{6.1 Aspects That Are Novel}\label{aspects-that-are-novel}

\begin{enumerate}
\def\labelenumi{\arabic{enumi}.}
\item
  \textbf{State-conditioned gatekeeper with Lyapunov convergence}: The
  combination of a discrete state partition (from clustering) with a
  Lyapunov-stabilized update rule is absent from all four compared
  frameworks.
\item
  \textbf{Multi-expert consensus as an acquisition signal}: BO uses GP
  uncertainty; AL uses model confidence; SCX uses multi-expert consensus
  \(C(x)\), which leverages \(M\) independent hypotheses rather than a
  single model's uncertainty.
\item
  \textbf{Memory-augmented self-correction with finite-time guarantee}:
  The monotonic memory bank + gatekeeper update loop with guaranteed
  finite-time termination (Theorem SE-2) has no direct analog in the
  compared frameworks.
\item
  \textbf{Completeness-incompleteness duality}: The simultaneous
  guarantee of termination (Theorem SE-2) and impossibility of
  self-certification (Claims SE-C1, SE-C2) is a structural feature not
  articulated in AL, BO, or AlphaZero.
\end{enumerate}

\subsubsection{6.2 Aspects That Are Not
Novel}\label{aspects-that-are-not-novel}

\begin{enumerate}
\def\labelenumi{\arabic{enumi}.}
\item
  \textbf{Sequential experimental design}: The iterative
  selection-evaluation-update loop is the defining characteristic of all
  four compared frameworks. SCX does not invent a new paradigm here ---
  it adapts existing paradigms to the NEP setting.
\item
  \textbf{Uncertainty-guided exploration}: Every framework in the
  comparison uses some form of uncertainty to guide exploration. SCX's
  specific mechanism (multi-expert consensus) is novel, but the
  principle is not.
\item
  \textbf{Thompson sampling / posterior sampling}: SCX's state-level
  value function \(V(s)\) is a form of posterior-sampling-based
  acquisition, conceptually similar to Thompson sampling in BO.
\item
  \textbf{Occam's razor via model selection}: The stability diagnostic
  (Proposition 6) implements a form of Occam's razor that is
  conceptually similar to Bayesian model selection, though the
  implementation is different.
\end{enumerate}

\subsubsection{6.3 The SCX Synthesis}\label{the-scx-synthesis}

The SCX self-evolution framework occupies a \textbf{unique position} in
the theory landscape:

\begin{itemize}
\tightlist
\item
  \textbf{Unlike AlphaZero}: It deals with real physical data (not
  synthetic), has unknown environment dynamics, and provides explicit
  convergence guarantees.
\item
  \textbf{Unlike BO}: It operates at the state level (not point-wise),
  handles multiple experts (not a single surrogate), and addresses the
  epistemic limitation of self-certification.
\item
  \textbf{Unlike AL}: It uses state-level certification (not point-wise
  queries), aggregates multiple independent experts (not a single
  model), and provides guarantees against label noise.
\item
  \textbf{Unlike Solomonoff induction}: It is computationally tractable,
  operates over a restricted hypothesis class, and converges in finite
  time --- at the cost of not guaranteeing convergence to truth.
\end{itemize}

The most accurate description is: \textbf{SCX is a restricted-domain,
computationally tractable approximation of Solomonoff induction,
operationalized through BO-inspired acquisition on a state space
discovered by AL-like clustering, stabilized by AlphaZero-inspired
monotonic improvement.}

\begin{center}\rule{0.5\linewidth}{0.5pt}\end{center}

\subsection{References}\label{references}

\begin{enumerate}
\def\labelenumi{\arabic{enumi}.}
\tightlist
\item
  Silver, D., et al.~(2017). Mastering Chess and Shogi by Self-Play with
  a General Reinforcement Learning Algorithm. \emph{arXiv:1712.01815}.
\item
  Silver, D., et al.~(2018). A general reinforcement learning algorithm
  that masters chess, shogi, and Go through self-play. \emph{Science},
  362(6419), 1140-1144.
\item
  Mockus, J. (1975). On Bayesian methods for seeking the extremum.
  \emph{IFIP Technical Conference}.
\item
  Jones, D. R., et al.~(1998). Efficient global optimization of
  expensive black-box functions. \emph{Journal of Global Optimization},
  13(4), 455-492.
\item
  Srinivas, N., et al.~(2010). Gaussian process optimization in the
  bandit setting: No regret and experimental design. \emph{ICML 2010}.
\item
  Settles, B. (2009). Active learning literature survey. \emph{Computer
  Sciences Technical Report 1648}, UW-Madison.
\item
  Balcan, M.-F., et al.~(2009). Agnostic active learning. \emph{Journal
  of Computer and System Sciences}, 75(1), 78-89.
\item
  Solomonoff, R. J. (1964). A formal theory of inductive inference. Part
  I. \emph{Information and Control}, 7(1), 1-22.
\item
  Hutter, M. (2005). \emph{Universal Artificial Intelligence}. Springer.
\item
  Li, M., \& Vitanyi, P. (2008). \emph{An Introduction to Kolmogorov
  Complexity and Its Applications} (3rd ed.). Springer.
\end{enumerate}

\begin{center}\rule{0.5\linewidth}{0.5pt}\end{center}

\emph{End of 08\_theory\_connections.md}

\end{document}
